\documentclass[12pt, letterpaper]{article}

%-------Packages---------
\usepackage{newpxtext,newpxmath} % Palatino-like text & math (modern)
\usepackage{bboldx}
\usepackage{mathtools}
\usepackage{tikz}
\usetikzlibrary{calc,intersections}
\usepackage{tikz-cd}
\usepackage[all,arc]{xy}
\usepackage{graphicx}

\usepackage{enumerate}
\usepackage[dvipsnames]{xcolor}
\usepackage{float}
\usepackage[left=1in,top=1in,right=1in,bottom=1in]{geometry}
\usepackage{fancyhdr}
\usepackage{multicol}
\usepackage{setspace}

\usepackage{dsfont}
\usepackage{mathrsfs}

\usepackage{tcolorbox}
\tcbuselibrary{theorems,skins,breakable}

\usepackage{xparse}
\usepackage{import}
\usepackage[utf8]{inputenc}

\usepackage{hyperref}
\usepackage{cleveref}

\usepackage{xcolor, ifthen, tcolorbox}
\tcbuselibrary{theorems,skins,breakable}
% Theorem System
% The following environments are provided:
%   New Problem:        \begin{problem} ...
%   New Question Header:   \qh (then followed by subproblems)
%   Subproblem:     \begin{subproblem} ...
%   Problem Proof:  \begin{problemproof} ...
%   Proof:          \\begin{pf}{color} ...
%   Remark:         \rmk (sentence), \rmkb (block)

% Define custom colors
\definecolor{thmscol}{HTML}{F4565B} %For Problems
\definecolor{lemscol}{HTML}{FE844C} %For Lemmes
\definecolor{prpscol}{HTML}{00A7EF} %For proposition
\definecolor{corscol}{HTML}{23DAFF} %For corrolaries
\definecolor{rmkscol}{HTML}{6DEEC5} %For remarks
\definecolor{defscol}{HTML}{18C991} %For definitions
\definecolor{exmscol}{HTML}{7DB800} %For examples
\definecolor{clmscol}{HTML}{165c58} %For claims

%Change between dark(black)/light(white) mode
\newcommand{\colormode}{white}
\ifthenelse{\equal{\colormode}{white}}
      {\newcommand{\TextColor}{black}}
      {\newcommand{\TextColor}{white}}
\newcommand{\pbcolormain}{thmscol!80!\colormode}
\newcommand{\pbcolorsecond}{thmscol!38!\colormode}


% ============================
% Problem Counter
% ============================
\newcounter{subproblemcounter}
\newcounter{problemcounter}

\newcommand{\q}{
    \setcounter{subproblemcounter}{0}%
    \stepcounter{problemcounter} % Increment problem counter
    \color{\TextColor}Problem \arabic{problemcounter}: % Display the problem counter
}

\newcommand{\stepsub}{
    \stepcounter{subproblemcounter}%
    \color{\TextColor}(\alph{subproblemcounter})
}
% ============================
% Problem
% ============================
\newtcolorbox{pb}[1]{
  enhanced,
  frame hidden,
  titlerule=0mm,
  toptitle=1mm,
  bottomtitle=1mm,
  fonttitle=\bfseries\large,
  coltitle=black,
  colbacktitle=\pbcolormain,
  colback=\pbcolorsecond,
  title={\q #1},
  coltext=\TextColor
}

\newtcolorbox{spb}[1]{
  enhanced,
  frame hidden,
  titlerule=0mm,
  toptitle=1mm,
  bottomtitle=1mm,
  fonttitle=\bfseries\large,
  coltitle=black,
  colbacktitle=\pbcolormain,
  colback=\pbcolorsecond,
  title={\stepsub #1},
  coltext=\TextColor,
  attach boxed title to top left={xshift=3mm,yshift=-2mm},
  boxed title style={
    colback=\pbcolormain,
    colframe=\pbcolormain,
  }
}

\newtcolorbox{pbh}[1]{
    enhanced,
    frame hidden,
    titlerule=0mm,
    toptitle=1mm,
    bottomtitle=1mm,
    fonttitle=\bfseries\large,
    coltitle=black,
    colbacktitle=\pbcolormain,
    colback=\pbcolorsecond,
    title={\q #1},
    boxrule=0pt,
    interior hidden,
    attach boxed title to top left={xshift=3mm,yshift=-2mm},
    boxed title style={
        colback=\pbcolormain,
        colframe=\pbcolormain,
    }
}

\newenvironment{problem}[1]{%
  \newpage
  \begin{pb}{#1}
    }{
  \end{pb}
}

\newenvironment{subproblem}[1]{%
  \begin{spb}{#1}
    }{
  \end{spb}
}

\newenvironment{problemproof}{%
    \begin{pf}{\pbcolormain}
        }{
    \end{pf}
}

\newcommand{\qh}[1][]{
    \newpage
    \begin{pbh}{#1}\end{pbh} % Display the problem counter
}

% ============================
% Claim
% ============================

\newtcbtheorem[number within=problemcounter]{claim}{Claim}
{
    enhanced,
    frame hidden,
    titlerule=0mm,
    toptitle=1mm,
    bottomtitle=1mm,
    fonttitle=\bfseries\large,
    coltitle=black,
    colbacktitle=clmscol!50!\colormode,
    colback=clmscol!20!\colormode,
}{clm}

\NewDocumentCommand{\clm}{m+m}{
    \begin{claim*}{#1}{}
        #2
    \end{claim*}
}

\newenvironment{clmpf}{
	{\noindent{\it \textbf{Proof.}}
	\setlength{\parindent}{0pt}
	}
	\tcolorbox[blanker,breakable,left=5mm,parbox=false,
    before upper={\parindent15pt},
    after skip=10pt,
	borderline west={1mm}{0pt}{clmscol!50!\colormode}]
}{
    \textcolor{clmscol!50!\colormode}{\hbox{}\nobreak\hfill$\blacksquare$}
    \endtcolorbox
}

\NewDocumentCommand{\clmp}{m+m+m}{
    \begin{claim*}{#1}{}
        #2
    \end{claim*}

    \begin{clmpf}
        #3
    \end{clmpf}
}


% ============================
% Proof
% ============================

\newenvironment{pf}[1]{%
  \def\pfcolor{#1}% Store the color in a macro
  \noindent{\it \textbf{Proof.}}%
  \tcolorbox[blanker,breakable,left=5mm,parbox=false,%
    before upper={\parindent15pt},%
    after skip=10pt,%
    borderline west={1mm}{0pt}{\pfcolor},%
    coltext=\TextColor
  ]%
  \setlength{\parindent}{0pt}%
}{%
  \textcolor{\pfcolor}{\hbox{}\nobreak\hfill$\blacksquare$}%
  \endtcolorbox%
}

\newtcolorbox{solution}{
  blanker,
  breakable,
  left=5mm,
  parbox=false,%
  before upper={\parindent15pt},%
  after skip=10pt,%
  borderline west={1mm}{0pt}{\pbcolormain},%
  coltext=\TextColor
}


% ============================
% Remark
% ============================


\NewDocumentCommand{\rmk}{+m}{
    {\it \color{rmkscol!100!white}#1}
}

\newenvironment{remark}{
    \par
    \vspace{5pt}
    \begin{minipage}{\textwidth}
        {\par\noindent{\textbf{Remark.}}}
        \tcolorbox[blanker,breakable,left=5mm,
        before skip=10pt,after skip=10pt,
        borderline west={1mm}{0pt}{rmkscol!80!\colormode},
        coltext=\TextColor
        ]
}{
        \endtcolorbox
    \end{minipage}
    \vspace{5pt}
}

\NewDocumentCommand{\rmkb}{+m}{
    \begin{remark}
        #1
    \end{remark}
}


% Define a new command for problems
\renewcommand{\headrule}{\color{\TextColor}\hrulefill}
\newcommand{\C}{\mathbb{C}}
\newcommand{\R}{\mathbb{R}}
\newcommand{\Q}{\mathbb{Q}}
\newcommand{\Z}{\mathbb{Z}}
\newcommand{\N}{\mathbb{N}}
\renewcommand{\P}{\mathbb{P}}
\newcommand{\F}{\mathbb{F}}
\newcommand{\contr}{\Rightarrow \Leftarrow}
\newcommand{\s}{\(\,\)}
\renewcommand{\t}[1]{\text{#1}}

\newcommand{\skcgen}{{\sf Gen}} 
\newcommand{\skcenc}{{\sf Enc}}
\newcommand{\skcdec}{{\sf Dec}}
\newcommand{\mac}{{\sf MAC}}
\newcommand{\ver}{{\sf Vrfy}}
\newcommand{\bit}{{\{0,1\}}}

% Meta data
\setstretch{1.2}

\begin{document}
\color{\TextColor}
\pagecolor{\colormode}
\setlength{\parindent}{0pt}
\pagestyle{fancy}
\fancyhf{}
\fancyhead[LOE]{\color{\TextColor}\slshape{\textbf{Vincent Ferrara}}}
\fancyhead[ROE]{\color{\TextColor}HW3 --- \thepage}
\fancyhead[C]{\color{\TextColor}EECS 475}

\begin{problem}{}
    Let $F\colon \bit^n \times \bit^n \rightarrow \bit^n$ be pseudorandom function.
  Show that each of the following MACs is insecure, even if used to
  authenticate fixed-length
  messages. (In each case $\skcgen$ outputs a uniform $k \in \bit^n$,
  $n$ being even, and $\langle i \rangle$ denotes an $n/2$-bit binary
  encoding of the integer $i$).
\end{problem}
\begin{subproblem}{}
    To authenticate a message $m = m_1\| \cdots \| m_\ell$, where
          $m_i \in \bit^n$, compute \\
          $t := F_k(m_1)\oplus \cdots \oplus F_k(m_\ell)$.
\end{subproblem}
\begin{problemproof}
    Let $m_1,m_2,m_3 \in \bit^n$ s.t. they are all distinct. Then query the oracle twice on
    $\mathcal{O}(m_1||m_2)=F_k(m_1) \oplus F_k(m_2)$ and $\mathcal{O}(m_1||m_3)=F_k(m_1) \oplus F_k(m_3)$.

    Then calculate $(F_k(m_1) \oplus F_k(m_2)) \oplus (F_k(m_1) \oplus F_k(m_3)) = F_k(m_2) \oplus F_k(m_3)=t$.

    Then, return $(m_2||m_3, t)$ which is a valid pair.
\end{problemproof}

\begin{subproblem}{}
    To authenticate a message $m = m_1 \| \cdots \| m_\ell$, where
          $m_i \in \bit^{n/2}$, compute \\
          $t := F_k(\langle 1 \rangle \| m_1) \oplus \cdots \oplus
            F_k(\langle \ell \rangle \| m_\ell)$.
\end{subproblem}
\begin{problemproof}
    Let $m_1,m_2,m_3 \in \bit^{n/2}$ s.t. they are all distinct. Then query the oracle three times on
    $\mathcal{O}(m_1||m_3)=F_k(\langle 1 \rangle \| m_1) \oplus F_k(\langle 2 \rangle \| m_3)$
    
    $\mathcal{O}(m_3 || m_2) = F_k(\langle 1 \rangle \| m_3) \oplus F_k(\langle 2 \rangle \| m_2)$.
    
    $\mathcal{O}(m_3 || m_3)= F_k(\langle 1 \rangle \| m_3) \oplus F_k(\langle 2 \rangle \| m_3)$.

    Then calculate $\mathcal{O}(m_1||m_3) \oplus \mathcal{O}(m_3 || m_2) \oplus \mathcal{O}(m_3 || m_3) = F_k(\langle 1 \rangle \| m_1) \oplus F_k(\langle 2 \rangle \| m_2)=t$.

    Then return $(m_1||m_2, t)$ which is a valid pair.
\end{problemproof}

\begin{problem}{}
    Prove that if a MAC scheme $(\mathsf{Gen}, \mathsf{Mac}, \mathsf{Vrfy})$ is strongly existentially unforgable under a chosen-message attack, then it is existentially unforgable under a chosen-message attack.
\end{problem}
\begin{problemproof}
    Assume that a MAC scheme $(\mathsf{Gen}, \mathsf{Mac}, \mathsf{Vrfy})$ is not existentially unforgable under a chosen-message attack.

    This implies there exists an adversary $\mathcal{A}$ can produce a valid pair of message and tag with non-negligible probability s.t. the message was never queried before.

    Define $C$ s.t.
    \[\begin{tikzcd}
	\begin{array}{c} \square\\\ k \leftarrow \skcgen(1^n) \end{array} && C && {\mathcal{A}} \\
	{} && {} && {}
	\arrow["{\mathsf{MAC}(m,k)=t}"', shift right=3, from=1-1, to=1-3]
	\arrow["m"', shift right=3, from=1-3, to=1-1]
	\arrow["t"', shift right=3, from=1-3, to=1-5]
	\arrow["m"', shift right=3, from=1-5, to=1-3]
	\arrow[shift right=5, from=1-5, to=1-5, loop, in=325, out=35, distance=10mm]
	\arrow["{(m',t')}", shift right=3, from=2-3, to=2-1]
	\arrow["{(m',t')}", shift right=3, from=2-5, to=2-3]
\end{tikzcd}\]
   Since $\mathcal{A}$ produces a valid pair with non-negligible probability s.t. $m'$ has not been queried before this implies that $C$ also produces a valid pair with non-negligible probability s.t. $(m',t')$ is a unique pair.
   This is a unique pair since $m'$ is unique and has not been seen before.

   Therefore, if a MAC scheme is strongly existentially unforgable under a chosen-message attack, then it is existentially unforgable under a chosen-message attack.
\end{problemproof}

\begin{problem}{}
    Let $H\colon \bit^\ell \to \bit^n$ be a collision-resistant hash function,
    where $\ell > n$, and $F\colon \mathcal{K} \times \bit^n \to \bit^n$ be a PRF.
    Show that the function $F'\colon \mathcal{K} \times \bit^\ell \to \bit^n$
    given by $F'_k(x) := F_k(H(x))$ is a PRF.  In practice, this construction
    can produce a PRF with arbitrary-length inputs.
\end{problem}
\begin{problemproof}
    Assume there exists a decider $D$ that can decide between a random function $R$ from $\ell$ bits to $n$ and $F \circ H$.

    Case 1: There exists queries from $D$ s.t. $x_i \neq x_j$ but $H(x_i)=H(x_j)$.

    Case 2: Is Case 1 not happening.

    Thus, 
    
    $\P[D \t{ decides correct}] = \P[D \t{ decides correct and Case 1}] + \P[D \t{ decides and Case 2}]$.

    Assume $\P[D \t{ decides correct and Case 1}]$ is non-negligible. Then define $C$ s.t.

    \[\begin{tikzcd}
	\square && \begin{array}{c} C\\ H^s(x) \end{array} && D
	\arrow["s"', from=1-1, to=1-3]
	\arrow["{y \overset{\$}{\leftarrow}U_n}"', shift right=3, from=1-3, to=1-5]
	\arrow["{x_i}"', shift right=3, from=1-5, to=1-3] \end{tikzcd}\]
    First get $x_i$ from $D$ and save $(x_i, H^s(x_i))$ then check
    if we got $H^s(x_i)$ before where the input is different if this happens we are done and return the pair we got.

    If we do not get a match then return $D$ a uniformly random $n$-bit string.

    Since the probability of getting a collision from $D$ is non-negligible this implies that $C$ also makes a collision of $H$ with non-negligible probability.

    Assume $\P[D \t{ decides correct and Case 2}]$ is non-negligible. Then define $C$ s.t.
    \[\begin{tikzcd}
	\begin{array}{c} \square\\\ k \leftarrow \mathcal{K} \end{array} && C && {D} \\
	{} && {} && {}
	\arrow["{\mathcal{O}(H(x))=y}"', shift right=3, from=1-1, to=1-3]
	\arrow["{H(x)}"', shift right=3, from=1-3, to=1-1]
	\arrow["y"', shift right=3, from=1-3, to=1-5]
	\arrow["x"', shift right=3, from=1-5, to=1-3]
	\arrow[shift right=5, from=1-5, to=1-5, loop, in=325, out=35, distance=10mm]
	\arrow["{b'}", shift right=3, from=2-3, to=2-1]
	\arrow["{b'}", shift right=3, from=2-5, to=2-3]
    \end{tikzcd}\]
    Since we are assuming that the collisions are negligible in the queries this implies that $H$ restricted to the set $Q = \{x \mid x \t{ was queried}\}$ is injective with 1 - negligible probability. This implies that

    Each $x \in Q$ maps to a unique hash value h$ = H(x)$.
    When we compose $R \circ H$, we are simply asking $R_0$ for independent uniform outputs on the distinct points $H(Q)\subseteq{0,1}^n$.
    For a random function $R$, the outputs $R(h)$ for distinct $h$ are independent and uniformly random in ${0,1}^n$.
    Thus, the induced mapping $x \mapsto R(H(x))$ assigns independent uniform outputs to each distinct $x\in Q$.
    This is exactly the defining property of a truly random function $R_\ell : {0,1}^\ell \to {0,1}^n$, since $R_\ell(x_i)$ are also independent and  uniformly random for distinct $x_i$.
    
    This implies that the induced map $R \circ H \mid_{Q}$ is indistinguishable to a random function $R_\ell$ from $\ell$ bits to $n$ bits.

    Therefore, since $D$ decides between $F'$ and $R_\ell \approx R \circ H$(given $H$ is injective on the queries) with non-negligible probability then $C$ will decide between $F$ and $R$ with non-negligible probability
\end{problemproof}
\begin{problem}{}
    Let $n\in\Z$ be a positive integer. Consider the
    following set \[\mathbb{G}_n := \{ an + 1 \mid a \in \{0, \ldots, n-1\} \}.\]
    Every element of $\mathbb{G}_n$ also belongs to the multiplicative group $\Z_{n^2}^*$ (recall that  $\Z_{n^2}^*$ contains all integers from $\{0,...,n^2-1\}$ that are coprime with $n^2$.)

    Show that $\mathbb{G}_n$ is a multiplicative subgroup of $\Z_{n^2}^*$.
\end{problem}
\begin{problemproof}
    (Associativity): Follows from multiplication.
    
    (Closure):

    Let $g,g' \in \mathbb{G}_n$ then $g=an+1$, and $g'=bn+1$ for $a,b \in \{0,\dots,n-1\}$.

    Consider $g \cdot g' = (an+1)(bn+1)=abn^2+an+bn+1=(a+b)n+1$ if $(a+b) < n$ then we are done. 

    If $(a+b) \geq n$ then $a+b=n+c$ for $c \in \{0,\dots,n-1\}$ since the max $(a+b)$ could be is $2n-2$, so we would get $n+n-2$, so $n-2$ is the max $c$ could be. Thus, $(a+b)n+1=(n+c)n+1 = cn+1 \in \mathbb{G}_n$. Thus, we have Closure.

    (Identity): Let $a = 0$ then $1 \in \mathbb{G}_n$ and $1 \cdot g = g \cdot 1 = g$ for all $g \in \mathbb{G}_n$.

    (Inverse): Let $g=an+1 \in \mathbb{G}_n$ for $a \in \{0,\dots,n-1\}$ then consider $g'=(n-a)n+1$.

    Thus, $g \cdot g' = (an+1)((n-a)n+1)=a(n-a)n^2+an+(n-a)n+1=an+n^2-an+1=1$. Also, $g' \cdot g = 1$.

    Thus, $\mathbb{G}_n$ is a subgroup of $\Z_{n^2}^*$
\end{problemproof}

\qh
\begin{subproblem}{}
    Show that if there exists a p.p.t. adversary $\mathcal{A}$ that breaks the DLOG problem, then there exists a p.p.t. adversary $\mathcal{B}$ that breaks the CDH problem.
\end{subproblem}
\begin{problemproof}
    Assume there exists a p.p.t. adversary $\mathcal{A}$ that breaks the DLOG problem.

    Define $C$ s.t. first get $(\mathbb{G},q,g,g^a,g^b)$ from the oracle. Then send $g^a$ to $\mathcal{A}$ and get back $a$ with non-negligible probability. Then send $g^b$ to $\mathcal{A}$  and get back $b$ with non-negligible probability. Then return $(g^a)^b$.

    Thus, $C$ breaks CDH problem with non-negligible probability.
\end{problemproof}

\begin{subproblem}{}
    Show that if there exists a p.p.t. adversary $\mathcal{B}$ that breaks the CDH problem, then there exists a p.p.t. adversary $\mathcal{C}$ that breaks the DDH problem.
\end{subproblem}
\begin{problemproof}
    Assume there exists a p.p.p adversary $\mathcal{B}$ that breaks the CHD problem.

    Define $C$ s.t. first get $(\mathbb{G},q,g^x,g^y,g^z)$ s.t. $z$ is either random or $xy$. Then send $g^a$ and $g^b$ to $\mathcal{B}$ and get back $g^{ab}$ with non-negligible probability. Then check whether $g^z=g^{ab}$. If this is true return 1 else return 0.

    Thus, $C$ breaks DDH problem with non-negligible probability.
\end{problemproof}
\end{document}